\section{The syntax of circuits on parties}
\label{sec:party_syntax}

Throughout this chapter $P$ is a finite set of parties, a system
$X=(X_p)_{p\in P}$ assigns a finite basis set to every party,
$\mathcal H_X=\bigotimes_p\C^{X_p}$, and $X_P$ and $X_D$ are the configurations of
all parties and of a set $D$ of parties, as in Chapter~\ref{ch:sites}. This
section defines the circuits run by the parties
\cite[Section~5.1]{OpenAI2026PolynomialPEPS}, reads them as circuits on the
geometry of all finite sets of parties, and counts their gates. It then
introduces the gates of the compression argument: pair sources and pair effects,
their monomials and weighted sums, and words of pair sources.

\subsection{Party circuits and their operators}
\label{sec:party_syntax_circuits}

\begin{definition}[Circuits on a finite set of parties]
    \label{def:party_circuit}
    \lean{QuantumCircuit.Party.PartyCircuit, QuantumCircuit.Party.PartyCircuit.toCircuit,
        QuantumCircuit.Party.PartyCircuit.op, QuantumCircuit.Party.PartyCircuit.IsAtom}
    \leanok
    \uses{def:circuits_gate_family, def:circuits_subsets_geometry, def:circuits_gate_at,
        def:circuits_of_cells, thm:circuits_composition, def:circuits_single,
        thm:circuits_single_op}
    Let $\mathcal G$ be a gate family on $P$. A party circuit from a system $X$ to a
    system $Y$ with gates in $\mathcal G$ is built by the following rules
    \cite[Section~5.1]{OpenAI2026PolynomialPEPS}:
    \begin{enumerate}
        \item the empty circuit from $X$ to $X$;
        \item one gate $g$ of $\mathcal G$ on a set $D\subseteq P$, its cell, from $X$
              to a system $Y$ with $Y_s=X_s$ for every $s\notin D$; the other parties
              are spectators, and the cell may be empty;
        \item one local operation at a party $p$: a matrix $A$ from $\C^{X_p}$ to
              $\C^{Y_p}$, where $Y_s=X_s$ for every $s\neq p$;
        \item local operations at every party at once: matrices $A_p$ from
              $\C^{X_p}$ to $\C^{Y_p}$ for all $p\in P$;
        \item the composition $c'\circ c$ of a party circuit $c$ from $X$ to $Y$ and a
              party circuit $c'$ from $Y$ to $Z$, first $c$ and then $c'$.
    \end{enumerate}
    The party circuits given by the first four rules are the party circuits of one
    step. The circuit $\widehat c$ of a party circuit $c$ is the circuit on the
    geometry of all finite sets of parties defined rule by rule: the empty circuit;
    the circuit of the layer of $g$ on $D$; the circuit of the layer of $A$ at $p$;
    the circuit of the layer with no cell and the one-member family $\{A_p\}$ at every
    party $p$; and the composition of $\widehat c$ and $\widehat{c'}$. When every gate
    of $\mathcal G$ has one outcome, $\widehat c$ has one branch
    (Theorem~\ref{thm:party_circuit_op}), and for an interpretation of $\mathcal G$ the
    operator $T_c$ of $c$, a matrix from $\mathcal H_X$ to $\mathcal H_Y$, is the
    Kraus operator of that branch.
\end{definition}

\begin{theorem}[The circuit and the operator of a party circuit]
    \label{thm:party_circuit_op}
    \lean{QuantumCircuit.Party.PartyCircuit.isUnbranched_toCircuit,
        QuantumCircuit.Party.PartyCircuit.isSequential_toCircuit,
        QuantumCircuit.Party.PartyCircuit.isSingle_toCircuit,
        QuantumCircuit.Party.PartyCircuit.op_done, QuantumCircuit.Party.PartyCircuit.op_comp,
        QuantumCircuit.Party.PartyCircuit.op_gate, QuantumCircuit.Party.PartyCircuit.op_localOp,
        QuantumCircuit.Party.PartyCircuit.op_localOps,
        QuantumCircuit.Party.PartyCircuit.op_induction}
    \leanok
    \uses{def:party_circuit, def:circuits_local, def:circuits_gate_at, def:circuits_single,
        def:sites_place_between, def:elementary_finite_tensor_product}
    Let $c$ be a party circuit from $X$ to $Y$. Its circuit $\widehat c$ is unbranched
    and sequential: every layer has at most one cell, and a layer with a cell is the
    identity at the parties outside it. If every gate has one outcome, $\widehat c$ is
    single. For an interpretation of a gate family with one outcome per gate, the
    operators are, rule by rule,
    \begin{align}
        T_{\mathrm{empty}} & =1,                           &
        T_{c'\circ c}      & =T_{c'}T_c,            \notag   \\
        T_g                & =M_g\otimes1,                 &
        T_A                & =A\otimes1,                   &
        T_{(A_p)_p}        & =\bigotimes_pA_p,    \notag
    \end{align}
    where $M_g$ is the matrix of $g$ at its outcome, and $M_g\otimes1$ and $A\otimes1$
    are the placements between $X$ and $Y$ on $D$ and on $\{p\}$, with $A$ read on the
    configurations of $\{p\}$. Consequently, let a property of matrices from
    $\mathcal H_X$ to $\mathcal H_Y$, for all systems $Y$, hold for the identity of
    $\mathcal H_X$ and pass from $M$ to $T_sM$ for every party circuit $s$ of one step.
    Then it holds for $T_c$ for every party circuit $c$ from $X$.
\end{theorem}
\begin{proof}\leanok
    \uses{lem:circuits_duality, lem:circuits_local_comp, def:circuits_gate_at,
        def:circuits_single, thm:circuits_single_op, lem:circuits_gate_at,
        thm:circuits_of_cells}
    By induction on $c$. A gate, a local operation at one party and local operations
    at every party give the circuit of one layer, which is unbranched
    (Lemma~\ref{lem:circuits_duality}). Each of these layers is one step, and the
    circuit of a layer of one step is sequential
    (Definition~\ref{def:circuits_gate_at}); the layer of a gate with one outcome, and
    a layer without cells whose onsite families have one member, give single circuits
    (Definition~\ref{def:circuits_single}). The three properties pass to compositions
    (Lemma~\ref{lem:circuits_local_comp} and Definitions~\ref{def:circuits_gate_at}
    and~\ref{def:circuits_single}). The operator of a single circuit of one layer is a
    Kraus operator of the layer (Theorem~\ref{thm:circuits_single_op}); for the layer
    of a gate it is $M_g\otimes1$ at the unique outcome, for the layer of $A$ at $p$ it
    is $A\otimes1$ (Lemma~\ref{lem:circuits_gate_at}), and for the layer without cells
    it is $\bigotimes_pA_p$ (Theorem~\ref{thm:circuits_of_cells}). The operator of a
    composition of single circuits is $T_{c'}T_c$ by
    Theorem~\ref{thm:circuits_single_op}. For the last statement, show by induction on
    a circuit $c$ from $Y$ to $Z$ that the property passes from $M$ to $T_cM$: a circuit
    of one step is the hypothesis, and $T_{c'\circ c}M=T_{c'}(T_cM)$. Apply this to
    $M=1$.
\end{proof}

\begin{definition}[Allowed and locally contractive party circuits]
    \label{def:party_allowed}
    \lean{QuantumCircuit.Party.PartyCircuit.IsAllowed,
        QuantumCircuit.Party.PartyCircuit.IsLocalContractive}
    \leanok
    \uses{def:party_circuit}
    Let $\mathcal G$ be a gate family with one outcome per gate, with an
    interpretation, and let $\mathcal A$ be a condition on the gates of $\mathcal G$. A
    party circuit is allowed for $\mathcal A$ when its local operations are
    contractions, $\lVert A\rVert\le1$ and $\lVert A_p\rVert\le1$, and every gate $g$
    satisfies $\mathcal A$ and has $\lVert M_g\rVert\le1$
    \cite[Section~5.1]{OpenAI2026PolynomialPEPS}. A party circuit is locally
    contractive when its local operations are contractions; its gates are not
    constrained.
\end{definition}

\begin{theorem}[An allowed party circuit is a contraction]
    \label{thm:party_op_contraction}
    \lean{QuantumCircuit.Party.PartyCircuit.norm_op_le_one,
        QuantumCircuit.Party.PartyCircuit.IsAllowed.forall_toCircuit,
        QuantumCircuit.Party.PartyCircuit.IsAllowed.isLocalContractive}
    \leanok
    \uses{def:party_allowed, def:circuits_local, def:kraus_family}
    Let $c$ be an allowed party circuit from $X$ to $Y$. Every gate and every onsite
    family of its circuit $\widehat c$ is trace-nonincreasing, $c$ is locally
    contractive, and $\lVert T_c\rVert\le1$
    \cite[Section~5.1]{OpenAI2026PolynomialPEPS}.
\end{theorem}
\begin{proof}\leanok
    \uses{thm:circuits_single_op, lem:circuits_gate_at, lem:circuits_local_gate_unitary,
        lem:kraus_family_channel, lem:circuits_local_comp, thm:party_circuit_op}
    By induction on $c$. A family with one member $V$ is trace-nonincreasing if and
    only if $\lVert V\rVert\le1$ (Theorem~\ref{thm:circuits_single_op}); in
    particular the family $\{1\}$ is, being unitary and hence trace-preserving
    (Lemmas~\ref{lem:circuits_local_gate_unitary} and~\ref{lem:kraus_family_channel}).
    Hence the layer of an allowed gate and the layer of a contraction at one party are
    trace-nonincreasing at every gate and onsite family
    (Lemma~\ref{lem:circuits_gate_at}), and so is the layer of contractions at every
    party. The property passes to compositions (Lemma~\ref{lem:circuits_local_comp}).
    Local contractivity is part of allowedness. Since $\widehat c$ is single
    (Theorem~\ref{thm:party_circuit_op}), Theorem~\ref{thm:circuits_single_op} gives
    $\lVert T_c\rVert\le1$.
\end{proof}

\subsection{Gates along a party circuit and their counts}
\label{sec:party_syntax_counts}

A gate whose cell consists of exactly one party is private. In the geometry of all
finite sets of parties a private cell costs $0$ and every other cell, the empty one
included, costs $1$ (Definition~\ref{def:circuits_subsets_geometry}).

\begin{definition}[Gate occurrences, gate counts and participation]
    \label{def:party_gate_counts}
    \lean{QuantumCircuit.GateOccurrence, QuantumCircuit.Party.PartyCircuit.gateLocations,
        QuantumCircuit.Party.PartyCircuit.participants, QuantumCircuit.Party.PartyCircuit.gates,
        QuantumCircuit.Party.PartyCircuit.gateCells, QuantumCircuit.Party.PartyCircuit.numGates,
        QuantumCircuit.Party.PartyCircuit.gateCount,
        QuantumCircuit.Party.PartyCircuit.participation,
        QuantumCircuit.Party.PartyCircuit.participationCount,
        QuantumCircuit.Party.PartyCircuit.ForallGate}
    \leanok
    \uses{def:party_circuit, def:circuits_subsets_geometry}
    Let $c$ be a party circuit. Its gate occurrences form a finite set, one for each
    gate, kept distinct even when two gates coincide: there is none for the empty
    circuit and for local operations, one for a gate, and the disjoint union of those
    of $c$ and $c'$ for $c'\circ c$. The participants of an occurrence are the parties
    of its cell \cite[Section~5.1]{OpenAI2026PolynomialPEPS}. The gates of $c$, each
    with its cell and its two systems, and their cells form lists in chronological
    order. Read along $c$, with the constructors of a composition adding:
    \begin{enumerate}
        \item the number of gates $n(c)$ counts every gate, private ones and those on
              the empty cell included;
        \item the gate count $\#(c)$ is the sum of the costs of the cells of the gates,
              the number of gates that are not private;
        \item the participation $\pi_p(c)$ of a party $p$ is the number of gates that
              are not private and whose cell contains $p$;
        \item the participation count $\nu_p(c)$ is the number of gates whose cell
              contains $p$, private ones included.
    \end{enumerate}
    A condition on gates holds at every gate of $c$ when every gate of $c$, with its
    cell and systems, satisfies it.
\end{definition}

\begin{theorem}[Relations among the counts]
    \label{thm:party_gate_count_relations}
    \lean{QuantumCircuit.Party.PartyCircuit.gateCells_eq_map_gates,
        QuantumCircuit.Party.PartyCircuit.length_gates,
        QuantumCircuit.Party.PartyCircuit.card_gateLocations,
        QuantumCircuit.Party.PartyCircuit.gateCount_le_numGates,
        QuantumCircuit.Party.PartyCircuit.gateCount_eq_numGates,
        QuantumCircuit.Party.PartyCircuit.participation_eq_participationCount,
        QuantumCircuit.Party.PartyCircuit.forall_participants_of_forallGate,
        QuantumCircuit.Party.PartyCircuit.natCard_participants_mem,
        QuantumCircuit.Party.PartyCircuit.gateCount_gate_empty}
    \leanok
    \uses{def:party_gate_counts}
    For every party circuit $c$, the list of cells is the list of the cells of the
    gates, and the length of the list of gates and the number of gate occurrences are
    both $n(c)$. Always $\#(c)\le n(c)$, and when no gate of $c$ is private,
    $\#(c)=n(c)$ and $\pi_p(c)=\nu_p(c)$ for every party $p$. A condition on cells that
    holds at every gate holds at the participants of every occurrence. The number of
    occurrences whose participants contain $p$ is $\nu_p(c)$
    \cite[Section~5.1]{OpenAI2026PolynomialPEPS}. One gate on the empty cell has gate
    count $1$.
\end{theorem}
\begin{proof}\leanok
    \uses{def:circuits_subsets_geometry}
    By induction on $c$, the constructors of a composition adding on both sides. For a
    gate on $D$ the cost of $D$ is at most $1$, it is $1$ when $|D|\neq1$, and it is
    positive exactly when $|D|\neq1$ (Definition~\ref{def:circuits_subsets_geometry}).
    The occurrences of a composition containing $p$ are the disjoint union of those of
    the two parts.
\end{proof}

\begin{theorem}[The counts of the circuit of a party circuit]
    \label{thm:party_counts_circuit}
    \lean{QuantumCircuit.Party.PartyCircuit.gateCount_toCircuit,
        QuantumCircuit.Party.PartyCircuit.participation_toCircuit}
    \leanok
    \uses{def:party_gate_counts, def:circuits_costs}
    If every gate has one outcome, then for every party circuit $c$ the gate count of
    its circuit $\widehat c$ is $\#(c)$, and the number of gates of $\widehat c$
    meeting a party $p$ is $\pi_p(c)$.
\end{theorem}
\begin{proof}\leanok
    \uses{thm:party_circuit_op, lem:circuits_gate_at, def:circuits_costs, lem:circuits_seq}
    By induction on $c$. The cost of the circuit of one layer is the cost of the
    layer. The layer of a gate on $D$ has gate count the cost of $D$ and meets $p$ once
    exactly when $p\in D$ and $D$ has positive cost, and the layer of a local operation
    costs nothing (Lemma~\ref{lem:circuits_gate_at}); a layer without cells costs
    nothing (Definition~\ref{def:circuits_costs}). For a composition, $\widehat c$ is
    unbranched with a branch (Theorem~\ref{thm:party_circuit_op}), so the costs of
    $\widehat c$ and $\widehat{c'}$ add (Lemma~\ref{lem:circuits_seq}).
\end{proof}

\subsection{Relabelling bases and replacing gates}
\label{sec:party_syntax_relabel}

\begin{definition}[Relabelling the bases and replacing the gates]
    \label{def:party_relabel}
    \lean{QuantumCircuit.Party.PartyCircuit.relabel,
        QuantumCircuit.Party.PartyCircuit.gates_relabel,
        QuantumCircuit.Party.PartyCircuit.gateCells_relabel,
        QuantumCircuit.Party.PartyCircuit.numGates_relabel,
        QuantumCircuit.Party.PartyCircuit.gateCount_relabel,
        QuantumCircuit.Party.PartyCircuit.mapGates,
        QuantumCircuit.Party.PartyCircuit.mapGatesLocationsEquiv}
    \leanok
    \uses{def:party_circuit, def:party_gate_counts, def:circuits_gate_morphism}
    For bijections $e_s\colon Y_s\to X_s$, the relabelling $R_e$ is the party circuit
    from $X$ to $Y$ of local operations at every party, at $s$ the matrix from
    $\C^{X_s}$ to $\C^{Y_s}$ whose entry at $(y,x)$ is $1$ if $e_s(y)=x$ and $0$
    otherwise. It has no gate and no cell, so $n(R_e)=\#(R_e)=0$. For a morphism
    $\varphi$ of gate families, the party circuit $\varphi_*c$ replaces every gate $g$
    of $c$ by $\varphi(g)$ and keeps every other rule; the gate occurrences of $c$ and
    of $\varphi_*c$ are identified rule by rule.
\end{definition}

\begin{theorem}[The operator of a relabelling]
    \label{thm:party_relabel_op}
    \lean{QuantumCircuit.Party.PartyCircuit.op_relabel,
        QuantumCircuit.Party.PartyCircuit.op_relabel_mul,
        QuantumCircuit.Party.PartyCircuit.op_relabel_symm,
        QuantumCircuit.Party.PartyCircuit.isAllowed_relabel}
    \leanok
    \uses{def:party_relabel, def:party_allowed}
    For a gate family with one outcome per gate and any interpretation, the operator
    of $R_e$ is the matrix of the bijection $y\mapsto(e_s(y_s))_s$ of configurations:
    $(T_{R_e})_{yx}=1$ if $e_s(y_s)=x_s$ for every $s$ and $0$ otherwise. Hence, for
    every matrix $M$ with rows indexed by $X_P$, the product $T_{R_e}M$ is $M$ with its
    rows reindexed, $(T_{R_e}M)_{y,w}=M_{(e_s(y_s))_s,w}$; no inverse appears. For
    bijections $e_s\colon X_s\to Y_s$, the operator of the relabelling along their
    inverses is the matrix of the inverse of the bijection of configurations. A
    relabelling is allowed for every condition on gates.
\end{theorem}
\begin{proof}\leanok
    \uses{thm:party_circuit_op, thm:sites_product_span, lem:matrix_l2_contractions}
    By Theorem~\ref{thm:party_circuit_op}, $T_{R_e}$ is the Kronecker product of the
    matrices of the $e_s$, whose entry at $(y,x)$ is $\prod_s[e_s(y_s)=x_s]$, the
    entry of the matrix of the bijection of configurations
    (Theorem~\ref{thm:sites_product_span}). Multiplying by the matrix of a bijection
    reindexes the rows. The matrix of a bijection is a contraction
    (Lemma~\ref{lem:matrix_l2_contractions}).
\end{proof}

\begin{theorem}[Replacing the gates]
    \label{thm:party_map_gates}
    \lean{QuantumCircuit.Party.PartyCircuit.op_mapGates,
        QuantumCircuit.Party.PartyCircuit.gates_mapGates,
        QuantumCircuit.Party.PartyCircuit.gateCells_mapGates,
        QuantumCircuit.Party.PartyCircuit.numGates_mapGates,
        QuantumCircuit.Party.PartyCircuit.gateCount_mapGates,
        QuantumCircuit.Party.PartyCircuit.participation_mapGates,
        QuantumCircuit.Party.PartyCircuit.participants_mapGatesLocationsEquiv,
        QuantumCircuit.Party.PartyCircuit.IsAllowed.mapGates}
    \leanok
    \uses{def:party_relabel, def:party_allowed}
    Let $\varphi$ be a morphism of gate families. The list of gates of $\varphi_*c$ is
    the list of gates of $c$ with every gate replaced by its image, cells and systems
    kept. Hence $\varphi_*c$ has the cells, the number of gates, the gate count and the
    participations of $c$, and corresponding occurrences have the same participants.
    Let both families have one outcome per gate and let $\varphi$ be compatible with
    interpretations of the two families. Then $T_{\varphi_*c}=T_c$. If moreover $c$ is
    allowed for a condition $\mathcal A$, and the image of every gate satisfying
    $\mathcal A$ satisfies a condition $\mathcal A_2$, then $\varphi_*c$ is allowed for
    $\mathcal A_2$.
\end{theorem}
\begin{proof}\leanok
    \uses{thm:party_circuit_op, def:circuits_single, def:circuits_gate_morphism}
    By induction on $c$. For a gate $g$ on $D$, compatibility and uniqueness of the
    outcome give $M_{\varphi(g)}=M_g$, so the placements on $D$ agree
    (Theorem~\ref{thm:party_circuit_op}) and the norms agree; local operations and
    compositions are unchanged.
\end{proof}

\subsection{Fresh pair registers}
\label{sec:party_syntax_add_pair}

\begin{definition}[Adding fresh factors at two parties]
    \label{def:party_add_pair}
    \lean{QuantumCircuit.Party.addPair, QuantumCircuit.Party.addPairLeft,
        QuantumCircuit.Party.addPairRight, QuantumCircuit.Party.addPair_of_ne,
        QuantumCircuit.Party.addPair_eq_of_notMem, QuantumCircuit.Party.addPairConfigEquiv,
        QuantumCircuit.Party.addPairConfigEquiv_apply_fst,
        QuantumCircuit.Party.addPairConfigEquiv_apply_snd_left,
        QuantumCircuit.Party.addPairConfigEquiv_apply_snd_right,
        QuantumCircuit.Party.addPairConfigEquiv_apply_snd_of_ne,
        QuantumCircuit.Party.addPairConfigEquiv_symm_apply_left,
        QuantumCircuit.Party.addPairConfigEquiv_symm_apply_right,
        QuantumCircuit.Party.addPairConfigEquiv_symm_apply_of_ne,
        QuantumCircuit.Party.pairFactor, QuantumCircuit.Party.pairFactor.unique,
        QuantumCircuit.Party.addPairSiteEquiv, QuantumCircuit.Party.card_pairFactor,
        QuantumCircuit.Party.card_addPair, QuantumCircuit.Party.pairFactorConfigEquiv,
        QuantumCircuit.Party.addPairMatrix}
    \leanok
    \uses{def:sites_config}
    Let $p\neq q$ be parties and $\alpha,\beta$ finite sets. The system
    $X^{pq}_{\alpha\beta}$ has the basis $\alpha\times X_p$ at $p$,
    $\beta\times X_q$ at $q$, and $X_s$ at every other party; it agrees with $X$
    outside $\{p,q\}$. Its configurations are pairs of a configuration of the fresh
    factors and a configuration of $X$: the bijection
    $X^{pq}_{\alpha\beta,P}\to(\alpha\times\beta)\times X_P$ sends $z$ to $((a,b),x)$,
    where $z_p=(a,x_p)$, $z_q=(b,x_q)$ and $z_s=x_s$ for $s\notin\{p,q\}$
    \cite[Section~5.1]{OpenAI2026PolynomialPEPS}. The fresh factors alone form the
    system $F^{pq}_{\alpha\beta}$ with the basis $\alpha$ at $p$, $\beta$ at $q$, and a
    single basis vector at every other party; its bases are single points when $\alpha$
    and $\beta$ are, its configurations are the pairs $\alpha\times\beta$, and
    \begin{align}
        |F^{pq}_{\alpha\beta,s}|
         & =\bigl(|\alpha|\text{ if }s=p\text{, else }1\bigr)
        \bigl(|\beta|\text{ if }s=q\text{, else }1\bigr). \notag
    \end{align}
    At every party $s$ the basis of $X^{pq}_{\alpha\beta}$ is
    $F^{pq}_{\alpha\beta,s}\times X_s$, so
    $|X^{pq}_{\alpha\beta,s}|=|F^{pq}_{\alpha\beta,s}|\,|X_s|$. For matrices $C_s$
    from $\C^{Y_s}$ to $\C^{Y'_s}$, the operation $C_s$ beside the fresh factors is the
    matrix $1\otimes C_s$ from $\C^{Y^{pq}_{\alpha\beta,s}}$ to
    $\C^{Y'^{pq}_{\alpha\beta,s}}$, read through these bases.
\end{definition}

\begin{theorem}[Party operations beside the fresh factors]
    \label{thm:party_add_pair_matrix}
    \lean{QuantumCircuit.Party.addPairConfigEquiv_eq_siteEquiv,
        QuantumCircuit.Party.norm_addPairMatrix_le,
        QuantumCircuit.Party.piKronecker_addPairMatrix}
    \leanok
    \uses{def:party_add_pair, def:elementary_finite_tensor_product}
    The bijection of configurations of $X^{pq}_{\alpha\beta}$ is read party by party:
    it is the composite of the bijections
    $X^{pq}_{\alpha\beta,s}\to F^{pq}_{\alpha\beta,s}\times X_s$ at every party, of the
    splitting of a family of pairs into a configuration of $F^{pq}_{\alpha\beta}$ and
    one of $X$, and of the identification of the configurations of
    $F^{pq}_{\alpha\beta}$ with $\alpha\times\beta$. For matrices $C_s$ as in Definition~\ref{def:party_add_pair},
    the operation beside the fresh factors has norm at most $\lVert C_s\rVert$, and
    their Kronecker product over all parties is
    \begin{align}
        \bigotimes_s(1\otimes C_s) & =1_{\alpha\times\beta}\otimes\bigotimes_sC_s, \notag
    \end{align}
    the right side read along the bijections of configurations of
    $Y^{pq}_{\alpha\beta}$ and $Y'^{pq}_{\alpha\beta}$.
\end{theorem}
\begin{proof}\leanok
    \uses{thm:detailed_balance_tensor_norm, thm:sites_product_span}
    The first statement is checked on the fresh parts and at every party separately.
    Reindexing does not increase the norm, and
    $\lVert1\otimes C\rVert\le\lVert C\rVert$
    (Theorem~\ref{thm:detailed_balance_tensor_norm}). For the Kronecker product,
    compute an entry with the first statement: the product over the parties of the
    entries of the identity on the fresh factors is the entry of the identity on the
    configurations of $F^{pq}_{\alpha\beta}$, that is of $1_{\alpha\times\beta}$, and
    the product of the entries of the $C_s$ is the entry of $\bigotimes_sC_s$
    (Theorem~\ref{thm:sites_product_span}).
\end{proof}

\subsection{Pair sources, pair effects and monomials}
\label{sec:party_syntax_pair_gates}

For a bijection $e\colon m\to\alpha\times n$ and $v\in\C^\alpha$, the preparation
$J_e(v)$ is the matrix from $\C^n$ to $\C^m$ with entry $v_a$ at $(y,x)$ when
$e(y)=(a,x)$ and $0$ otherwise; it maps $\ket x$ to $v\otimes\ket x$ read through $e$.
The effect $\bra v\otimes1$ from $\C^{\alpha\times n}$ to $\C^n$ is its adjoint for
$e$ the identity, and $J_e(v)^\dagger$ is this effect with its columns read through $e$ (Definitions~\ref{def:construction_idle_gates}
and~\ref{def:party_effect_matrices}).

\begin{definition}[Pair sources, pair effects and monomials]
    \label{def:party_pair_gates}
    \lean{QuantumCircuit.Party.PairSource, QuantumCircuit.Party.SourceInventory,
        QuantumCircuit.Party.SourceInventory.IsNormalized,
        QuantumCircuit.Party.PairSource.partyPair,
        QuantumCircuit.Party.SourceGate, QuantumCircuit.Party.SourceGate.left_mem,
        QuantumCircuit.Party.SourceGate.right_mem, QuantumCircuit.Party.SourceGate.ofAddPair,
        QuantumCircuit.Party.SourceGate.regionMatrix,
        QuantumCircuit.Party.SourceGate.outputEquiv, QuantumCircuit.Party.SourceGate.ev,
        QuantumCircuit.Party.SourceGate.IsAllowed,
        QuantumCircuit.Party.SourceGate.ev_ofAddPair_apply,
        QuantumCircuit.Party.SourceGate.reindexOutput,
        QuantumCircuit.Party.SourceGate.ev_reindexOutput,
        QuantumCircuit.Party.SourceGate.ne_left_of_notMem,
        QuantumCircuit.Party.SourceGate.ne_right_of_notMem,
        QuantumCircuit.Party.SourceGate.toAddPair,
        QuantumCircuit.Party.SourceGate.addPairLeft_toAddPair,
        QuantumCircuit.Party.SourceGate.addPairRight_toAddPair,
        QuantumCircuit.Party.SourceGate.toAddPair_of_notMem,
        QuantumCircuit.Party.SourceGate.configEquiv,
        QuantumCircuit.Party.SourceGate.configEquiv_apply_fst,
        QuantumCircuit.Party.SourceGate.configEquiv_apply_snd_left,
        QuantumCircuit.Party.SourceGate.configEquiv_apply_snd_right,
        QuantumCircuit.Party.SourceGate.configEquiv_apply_snd_of_notMem,
        QuantumCircuit.Party.PairEffectGate, QuantumCircuit.Party.PairEffectGate.left_mem,
        QuantumCircuit.Party.PairEffectGate.right_mem,
        QuantumCircuit.Party.PairEffectGate.ofAddPair,
        QuantumCircuit.Party.PairEffectGate.toSourceGate,
        QuantumCircuit.Party.PairEffectGate.regionMatrix,
        QuantumCircuit.Party.PairEffectGate.reindexInput,
        QuantumCircuit.Party.PairEffectGate.ev, QuantumCircuit.Party.PairEffectGate.IsAllowed,
        QuantumCircuit.Party.PairEffectGate.ev_ofAddPair_apply,
        QuantumCircuit.Party.PairEffectGate.ev_reindexInput,
        QuantumCircuit.Party.PairEffectGate.toAddPair,
        QuantumCircuit.Party.PairEffectGate.configEquiv,
        QuantumCircuit.Party.ChainGate, QuantumCircuit.Party.chainEv,
        QuantumCircuit.Party.ChainGate.IsAllowed, QuantumCircuit.Party.Monomial}
    \leanok
    \uses{def:party_add_pair, def:party_circuit, def:circuits_gate_family,
        def:circuits_gate_sum, def:sites_pair_factor_config}
    A pair source datum consists of two distinct parties $l\neq r$, finite sets $\alpha$
    and $\beta$, and a vector $\eta\in\C^{\alpha\times\beta}$; its unordered pair of
    endpoints is $\{l,r\}$. A source inventory is a finite list of such data, repeated
    data counting as distinct occurrences, and it is normalized when every vector has
    norm $1$ \cite[Lemma~5.1]{OpenAI2026PolynomialPEPS}.

    A pair source on a cell $D$ from $X$ to $Y$ is a pair source datum with
    $D=\{l,r\}$, bijections $Y_l\to\alpha\times X_l$ and $Y_r\to\beta\times X_r$, and
    $Y_s=X_s$ for every $s\notin D$; thus $l,r\in D$, and a party outside $D$ is neither
    endpoint \cite[Section~5.1]{OpenAI2026PolynomialPEPS}. Its matrix on the cell, from
    $\C^{X_D}$ to $\C^{Y_D}$, prepares $\eta$ on the fresh factors and is the identity
    on the old ones: its entry at $(y,x)$ is $\eta_{ab}$ when $y_l\mapsto(a,x_l)$ and
    $y_r\mapsto(b,x_r)$, and $0$ otherwise. Correspondingly, a configuration of $D$ in
    $Y$ is a pair of a configuration of the fresh factors and a configuration of $D$ in
    $X$. A pair source has one outcome, its interpretation is its matrix, and it is
    allowed when $\lVert\eta\rVert=1$. The canonical source of $\eta$ on $\{p,q\}$ goes
    from $X$ to $X^{pq}_{\alpha\beta}$ with the identifications of
    Definition~\ref{def:party_add_pair}. Relabelling the output along bijections
    $e_s\colon Y'_s\to Y_s$, with $Y'_s=Y_s$ outside $D$, composes the two bijections of
    the source with $e_l$ and $e_r$ and reindexes the rows of the matrix by
    $y'\mapsto(e_s(y'_s))_s$. The output of a source is read in canonical form by the
    bijections $Y_s\to X^{lr}_{\alpha\beta,s}$ given at $l$ and $r$ by the bijections of
    the source and elsewhere by the identification $Y_s=X_s$; through them a
    configuration $y$ of $Y$ is the pair $((a,b),x)$ with
    $y_l\mapsto(a,x_l)$, $y_r\mapsto(b,x_r)$ and $x_s=y_s$ outside $D$.

    A pair effect on $D$ from $X$ to $Y$ has the same data with bijections
    $X_l\to\alpha\times Y_l$ and $X_r\to\beta\times Y_r$, and $Y_s=X_s$ outside $D$; it
    removes the factors $\alpha$ and $\beta$. Its matrix on the cell has the entry
    $\overline{\eta_{ab}}$ at $(y,x)$ when $x_l\mapsto(a,y_l)$ and $x_r\mapsto(b,y_r)$,
    and $0$ otherwise. The source with the same data goes from $Y$ to $X$. The
    canonical effect of $\eta$ on $\{p,q\}$ goes from $Y^{pq}_{\alpha\beta}$ to $Y$, a
    relabelling of the input reindexes the columns, and the input is read in canonical
    form through the source with the same data. A pair effect has one outcome, its
    interpretation is its matrix, and it is allowed when $\lVert\eta\rVert=1$
    \cite[Section~5.1]{OpenAI2026PolynomialPEPS}.

    The gates of monomials are the pair sources and the pair effects, with the
    interpretation of the sum of the two gate families; each has one outcome, and each
    is allowed when its vector is normalized. A monomial on a finite set $D$ of parties,
    from $X$ to $Y$ with $X,Y$ systems on $D$, is a party circuit on the parties of $D$
    whose gates are pair sources and pair effects.
\end{definition}

\begin{theorem}[Matrices of pair sources and pair effects]
    \label{thm:party_pair_gate_matrices}
    \lean{QuantumCircuit.Party.SourceGate.regionMatrix_eq_prepMatrix,
        QuantumCircuit.Party.PairEffectGate.regionMatrix_eq_conjTranspose,
        QuantumCircuit.Party.SourceGate.norm_ev_le, QuantumCircuit.Party.SourceGate.norm_ev,
        QuantumCircuit.Party.PairEffectGate.norm_ev_le,
        QuantumCircuit.Party.PairEffectGate.norm_ev,
        QuantumCircuit.Party.SourceGate.eq_reindexOutput,
        QuantumCircuit.Party.SourceGate.eq_reindexOutput_ofAddPair}
    \leanok
    \uses{def:party_pair_gates, def:party_effect_matrices}
    Let $S$ be the matrix of a pair source on $D$ with vector $\eta$. Then $S$ is the
    preparation $J_e(\eta)$ along the bijection $e\colon Y_D\to(\alpha\times\beta)\times X_D$,
    and the matrix of a pair effect is the adjoint of the matrix of the source with the
    same data. Both have norm at most $\lVert\eta\rVert$, with equality when $D$ has a
    configuration in the input of the source, respectively in the output of the effect;
    a normalized source or effect is a contraction. Two pair sources on $D$ from $X$
    with the same pair source datum differ by a relabelling of the output, and every
    pair source is, with $D=\{l,r\}$, the canonical source of its datum followed by a
    relabelling of the output.
\end{theorem}
\begin{proof}\leanok
    \uses{def:sites_pair_factor_config, def:party_effect_matrices}
    The entries of $S$ and of the preparation agree once the configurations of
    $D=\{l,r\}$ are read as pairs (Definition~\ref{def:sites_pair_factor_config}), and
    the entries of the effect are the conjugates of those of the source, transposed.
    The norms are those of a preparation (Definition~\ref{def:party_effect_matrices}),
    and $\lVert M^\dagger\rVert=\lVert M\rVert$. For two sources with bijections
    $u_l,u_r$ and $u'_l,u'_r$, the relabelling is $u_l^{-1}u'_l$ at $l$,
    $u_r^{-1}u'_r$ at $r$, and the identification of equal bases elsewhere; the
    canonical source is one such source.
\end{proof}

\begin{theorem}[The operator of one pair source or pair effect]
    \label{thm:party_pair_gate_op}
    \lean{QuantumCircuit.Party.SourceGate.configEquiv_ofAddPair,
        QuantumCircuit.Party.SourceGate.op_gate, QuantumCircuit.Party.PairEffectGate.op_gate,
        QuantumCircuit.Party.PairEffectGate.op_gate_eq_effectMatrix,
        QuantumCircuit.Party.ChainGate.op_gate_inl, QuantumCircuit.Party.ChainGate.op_gate_inr}
    \leanok
    \uses{def:party_pair_gates, def:party_effect_matrices, def:party_circuit}
    Let $g$ be a pair source on $D$ from $X$ to $Y$ with vector $\eta$, and let
    $\kappa\colon Y_P\to(\alpha\times\beta)\times X_P$ be the reading of its output in
    canonical form; for the canonical source $\kappa$ is the bijection of
    configurations of $X^{pq}_{\alpha\beta}$. The party circuit of the one gate $g$ has
    the operator $J_\kappa(\eta)$, which maps $\ket x$ to $\eta\otimes\ket x$.
    For a pair effect from $X$ to $Y$ with input read through
    $\kappa\colon X_P\to(\alpha\times\beta)\times Y_P$, the operator of its one gate is
    $J_\kappa(\eta)^\dagger$, the effect $\bra\eta\otimes1$ with its columns read
    through $\kappa$. Among the gates of monomials, one pair source and one pair effect
    have the same operators.
\end{theorem}
\begin{proof}\leanok
    \uses{thm:party_circuit_op, def:sites_place_between, lem:sites_place_between_algebra,
        thm:party_pair_gate_matrices, def:party_effect_matrices, def:circuits_gate_sum}
    The canonical reading is checked coordinate by coordinate. By
    Theorem~\ref{thm:party_circuit_op} the operator of the source is the placement
    $S\otimes1$ between $X$ and $Y$ on $D$, whose entry at $(y,x)$ is the entry of $S$
    when $y$ and $x$ agree outside $D$ (Definition~\ref{def:sites_place_between}). The
    two conditions, $y_l\mapsto(a,x_l)$, $y_r\mapsto(b,x_r)$, and agreement outside
    $D$, together say $\kappa(y)=((a,b),x)$, and the entry is then $\eta_{ab}$. For the
    effect, $(E\otimes1)=(S'^\dagger\otimes1)=(S'\otimes1)^\dagger$ for the source $S'$
    with the same data (Theorem~\ref{thm:party_pair_gate_matrices} and
    Lemma~\ref{lem:sites_place_between_algebra}), and the adjoint of a preparation is
    the effect with its columns reindexed (Definition~\ref{def:party_effect_matrices}).
    The interpretation of the sum of the two gate families restricts to the
    interpretations of each.
\end{proof}

\subsection{Weighted gates and original circuits}
\label{sec:party_syntax_original}

\begin{definition}[Weighted gates and effect circuits]
    \label{def:party_weighted_gate}
    \lean{QuantumCircuit.Party.WeightedGate, QuantumCircuit.Party.WeightedGate.ev,
        QuantumCircuit.Party.WeightedGate.IsAllowed, QuantumCircuit.Party.EffectCircuit}
    \leanok
    \uses{def:party_pair_gates, def:party_allowed}
    A weighted gate on a cell $D$ from $X$ to $Y$, its gate expansion, is a finite list
    of pairs $(\lambda_k,M_k)$ of a coefficient $\lambda_k\in\C$ and a monomial $M_k$ on
    the parties of $D$ from $X|_D$ to $Y|_D$
    \cite[Section~5.1]{OpenAI2026PolynomialPEPS}. It has one outcome, and its matrix,
    from $\C^{X_D}$ to $\C^{Y_D}$, is $\sum_k\lambda_kT_{M_k}$. It is allowed when every
    monomial $M_k$ is allowed: its local operations are contractions and its pair
    sources and pair effects are normalized. An effect circuit is a party circuit whose
    gates are weighted gates, with no further condition.
\end{definition}

\begin{definition}[Original gates and original circuits]
    \label{def:party_original_circuit}
    \lean{QuantumCircuit.Party.OriginalGate, QuantumCircuit.Party.OriginalGate.ev,
        QuantumCircuit.Party.OriginalGate.toWeightedGateMorphism,
        QuantumCircuit.Party.OriginalGate.isCompatible_toWeightedGateMorphism,
        QuantumCircuit.Party.OriginalGate.ofWeightedGate,
        QuantumCircuit.Party.OriginalGate.toWeightedGate_ofWeightedGate,
        QuantumCircuit.Party.OriginalCircuit, QuantumCircuit.Party.OriginalCircuit.done,
        QuantumCircuit.Party.OriginalCircuit.comp, QuantumCircuit.Party.OriginalCircuit.gate,
        QuantumCircuit.Party.OriginalCircuit.localOp,
        QuantumCircuit.Party.OriginalCircuit.localOps,
        QuantumCircuit.Party.OriginalCircuit.toPartyCircuit_done,
        QuantumCircuit.Party.OriginalCircuit.toPartyCircuit_comp,
        QuantumCircuit.Party.OriginalCircuit.toPartyCircuit_gate,
        QuantumCircuit.Party.OriginalCircuit.toPartyCircuit_localOp,
        QuantumCircuit.Party.OriginalCircuit.toPartyCircuit_localOps,
        QuantumCircuit.Party.OriginalCircuit.op, QuantumCircuit.Party.OriginalCircuit.op_done,
        QuantumCircuit.Party.OriginalCircuit.toEffectCircuit}
    \leanok
    \uses{def:party_weighted_gate, def:party_relabel, def:circuits_gate_morphism}
    An original gate on a cell $D$ from $X$ to $Y$ is a weighted gate on $D$ that is not
    private, $|D|\neq1$, whose monomials are allowed and whose matrix is a contraction;
    the empty cell is allowed \cite[Section~5.1]{OpenAI2026PolynomialPEPS}. It has one
    outcome and its matrix is that of its expansion. Forgetting the conditions is a
    morphism from original gates to weighted gates compatible with the
    interpretations, and a weighted gate satisfying the three conditions is the
    expansion of an original gate. An original circuit from $X$ to $Y$ is a party
    circuit whose gates are original gates and whose local operations are
    contractions. Original circuits are built by the empty circuit, one original gate,
    one contraction at one party, contractions at every party, and composition, and
    these are the corresponding party circuits. The operator $T_w$ of an original
    circuit $w$ is that of its party circuit, the empty circuit having the operator
    $1$, and its effect circuit is the image of its party circuit under the forgetful
    morphism.
\end{definition}

\begin{theorem}[Operators of original circuits]
    \label{thm:party_original_op}
    \lean{QuantumCircuit.Party.OriginalCircuit.induction,
        QuantumCircuit.Party.OriginalCircuit.op_comp,
        QuantumCircuit.Party.OriginalCircuit.op_gate,
        QuantumCircuit.Party.OriginalCircuit.op_localOp,
        QuantumCircuit.Party.OriginalCircuit.op_localOps,
        QuantumCircuit.Party.PartyCircuit.isAllowed_of_isLocalContractive,
        QuantumCircuit.Party.OriginalCircuit.isAllowed,
        QuantumCircuit.Party.OriginalCircuit.norm_op_le_one}
    \leanok
    \uses{def:party_original_circuit, def:party_allowed, def:sites_place_between,
        def:elementary_finite_tensor_product}
    A property of original circuits holds for all of them when it holds for the empty
    circuits, for one original gate, for one contraction at one party and for
    contractions at every party, and passes to compositions. The operator of a
    composition is $T_{w'\circ w}=T_{w'}T_w$, that of one original gate is the
    placement between $X$ and $Y$ on its cell of the matrix of its expansion, that of
    one contraction at a party $p$ is its placement on $\{p\}$, and that of
    contractions at every party is their Kronecker product. A locally contractive
    party circuit of original gates is allowed for the condition that the monomials of
    every gate are allowed; hence every original circuit is allowed, and
    $\lVert T_w\rVert\le1$ \cite[Section~5.1]{OpenAI2026PolynomialPEPS}.
\end{theorem}
\begin{proof}\leanok
    \uses{thm:party_circuit_op, thm:party_op_contraction}
    The induction is the induction on the underlying party circuit, the contractivity
    of the local operations of a composition splitting into that of the two parts. The
    operators are those of Theorem~\ref{thm:party_circuit_op}. Allowedness is checked
    by induction on the party circuit: every original gate is a contraction with
    allowed monomials. Theorem~\ref{thm:party_op_contraction} gives the norm bound.
\end{proof}

\begin{theorem}[The effect circuit of an original circuit]
    \label{thm:party_original_effect}
    \lean{QuantumCircuit.Party.PartyCircuit.forallGate_card_ne_one,
        QuantumCircuit.Party.PartyCircuit.forallGate_card_ne_one_mapGates,
        QuantumCircuit.Party.OriginalCircuit.forallGate_card_ne_one,
        QuantumCircuit.Party.OriginalCircuit.op_toEffectCircuit,
        QuantumCircuit.Party.OriginalCircuit.isAllowed_toEffectCircuit,
        QuantumCircuit.Party.OriginalCircuit.nonprivate}
    \leanok
    \uses{def:party_original_circuit, def:party_gate_counts, def:party_allowed}
    No gate of a party circuit of original gates is private, and neither is any gate of
    its image under the forgetful morphism; in particular no gate of an original
    circuit or of its effect circuit is private
    \cite[Section~5.1]{OpenAI2026PolynomialPEPS}. The effect circuit of an original
    circuit $w$ has the operator $T_w$ and is allowed for the condition that the
    monomials of every weighted gate are allowed.
\end{theorem}
\begin{proof}\leanok
    \uses{thm:party_map_gates, thm:party_original_op}
    The first statements follow by induction on the circuit from $|D|\neq1$ for every
    original gate, the forgetful morphism keeping the cells. The forgetful morphism is
    compatible with the interpretations and keeps the condition on the monomials, so
    Theorem~\ref{thm:party_map_gates} transfers the operator and allowedness of
    Theorem~\ref{thm:party_original_op}.
\end{proof}

\subsection{Words of pair sources}
\label{sec:party_syntax_words}

\begin{definition}[Words of pair sources]
    \label{def:party_word}
    \lean{QuantumCircuit.Party.Word, QuantumCircuit.Party.Word.IsAllowed,
        QuantumCircuit.Party.Word.op, QuantumCircuit.Party.IsSourceOnly,
        QuantumCircuit.Party.Word.source, QuantumCircuit.Party.Word.sourceAt,
        QuantumCircuit.Party.Word.localPrepare, QuantumCircuit.Party.Word.localDiscard,
        QuantumCircuit.Party.Word.prepareProduct, QuantumCircuit.Party.Word.sources,
        QuantumCircuit.Party.Word.sourceCount, QuantumCircuit.Party.Word.sources_comp,
        QuantumCircuit.Party.Word.sources_relabel, QuantumCircuit.Party.Word.sources_source,
        QuantumCircuit.Party.Word.sources_prepareProduct, QuantumCircuit.Party.Word.UsesOnly,
        QuantumCircuit.Party.Word.toChain}
    \leanok
    \uses{def:party_pair_gates, def:party_allowed, def:party_relabel,
        def:circuits_gate_sum}
    A word from $X$ to $Y$ is a party circuit whose gates are pair sources. It is
    allowed when its local operations are contractions and its pair sources are
    normalized, and an operator is source-only when it is the operator $T_w$ of an
    allowed word $w$ \cite[Lemma~5.1]{OpenAI2026PolynomialPEPS}. Particular words are:
    \begin{enumerate}
        \item the word of one pair source, and the canonical pair source word of
              $\eta\in\C^{\alpha\times\beta}$ on $p\neq q$, from $X$ to
              $X^{pq}_{\alpha\beta}$;
        \item the local preparation of $v\in\C^\alpha$ at a party $p$, the local
              operation $J_e(v)$ at $p$ for a bijection
              $e\colon Y_p\to\alpha\times X_p$, where $Y_s=X_s$ for $s\neq p$; with
              $\alpha=\alpha_1\times\alpha_2$ it prepares a pair vector at one party;
        \item the local bra of $v\in\C^\alpha$ at $p$, the local operation
              $J_e(v)^\dagger$ for a bijection $e\colon X_p\to\alpha\times Y_p$,
              where $Y_s=X_s$ for $s\neq p$;
        \item the preparation of the product vector $\bigotimes_p\phi_p$, for vectors
              $\phi_p\in\C^{X_p}$, from the system with a single basis vector at every
              party, by the local operation $\ket{\phi_p}$ at every party $p$.
    \end{enumerate}
    The source list of a word is the list of the data of its pair sources, those of the
    later part first: it is empty for the empty word and for local operations, it is
    $(g)$ for one source $g$, and it is the list of $w'$ followed by that of $w$ for
    $w'\circ w$; thus relabellings and product preparations have no source. The source
    count is the length of the source list. A word uses only the parties of a set $R$
    when its sources have cells inside $R$, its local operations at one party act at
    parties of $R$, and at every party $s\notin R$ its local operations at every party
    have $Y_s=X_s$ and are the identity. Every word is a monomial with no pair effect,
    each pair source read as a gate of monomials.
\end{definition}

\begin{theorem}[Operators of words]
    \label{thm:party_word_op}
    \lean{QuantumCircuit.Party.Word.op_comp, QuantumCircuit.Party.Word.op_source,
        QuantumCircuit.Party.Word.op_sourceAt,
        QuantumCircuit.Party.Word.op_source_eq_op_relabel_mul,
        QuantumCircuit.Party.Word.op_localPrepare,
        QuantumCircuit.Party.Word.op_localPrepare_eq_prepMatrix,
        QuantumCircuit.Party.Word.op_localDiscard,
        QuantumCircuit.Party.Word.op_prepareProduct_mulVec,
        QuantumCircuit.Party.Word.op_toChain}
    \leanok
    \uses{def:party_word, def:sites_pair_factor_config, def:sites_place_between,
        def:siteconfig_product_vector, def:party_effect_matrices}
    The operator of a composition of words is $T_{w'\circ w}=T_{w'}T_w$. The operator of
    one pair source is the preparation of its vector with the output read in canonical
    form (Theorem~\ref{thm:party_pair_gate_op}); for the canonical source on $p\neq q$ it
    is $J_e(\eta)$ along the bijection $e$ of configurations of
    $X^{pq}_{\alpha\beta}$, which maps $\ket x$ to $\eta\otimes\ket x$, and every pair
    source is the canonical source of its datum followed by the relabelling of its
    output into canonical form:
    \begin{align}
        T_g & =T_{R}\,T_{\mathrm{can}}, \notag
    \end{align}
    where $T_{\mathrm{can}}$ is the operator of the canonical source of the datum of $g$
    from $X$ to $X^{lr}_{\alpha\beta}$ and $R$ is the relabelling along the reading of
    the output of $g$ in canonical form.
    The local preparation of $v$ at $p$ has the operator $J_e(v)\otimes1$, the
    placement on $\{p\}$, which is the preparation $J_{\hat e}(v)$ along the
    bijection $\hat e\colon Y_P\to\alpha\times X_P$ that splits off the factor $\alpha$
    at $p$ (Definition~\ref{def:sites_pair_factor_config}). The local bra of $v$ at $p$
    has the operator $\bra v\otimes1$ with its columns read along the corresponding
    bijection $X_P\to\alpha\times Y_P$. The product preparation maps the unique basis
    vector to $\bigotimes_p\phi_p$. A word read as a monomial has the operator of the
    word.
\end{theorem}
\begin{proof}\leanok
    \uses{thm:party_circuit_op, thm:party_pair_gate_op, thm:party_relabel_op,
        def:sites_pair_factor_config, def:sites_place_between, thm:party_map_gates,
        def:circuits_gate_sum}
    The composition and the local operations are Theorem~\ref{thm:party_circuit_op},
    and one source is Theorem~\ref{thm:party_pair_gate_op}. The relabelling $R$ after
    the canonical source reindexes the rows of $J_e(\eta)$ along the reading in
    canonical form (Theorem~\ref{thm:party_relabel_op}), which is the operator of the
    source. For the local preparation compare entries: the entry of the placement at
    $(y,z)$ is $v_a$ when $e(y_p)=(a,z_p)$ and $y_s=z_s$ for $s\neq p$, which is the
    condition that $\hat e(y)=(a,z)$ (Definition~\ref{def:sites_pair_factor_config});
    the local bra is the same computation for the adjoint. For the product
    preparation, the Kronecker product of the columns $\ket{\phi_p}$ applied to the
    basis vector has the coefficient $\prod_p(\phi_p)_{x_p}$ at $x$. The inclusion of
    pair sources among the gates of monomials is compatible with the interpretations
    (Definition~\ref{def:circuits_gate_sum}), so Theorem~\ref{thm:party_map_gates}
    applies.
\end{proof}

\begin{theorem}[A fresh source commutes with local operations on the old factors]
    \label{thm:party_source_commutes}
    \lean{QuantumCircuit.Party.Word.op_sourceAt_mul_piKronecker}
    \leanok
    \uses{def:party_word, thm:party_add_pair_matrix}
    Let $p\neq q$, let $\eta\in\C^{\alpha\times\beta}$, and let $A_s$ be matrices from
    $\C^{Y_s}$ to $\C^{Y'_s}$. Local operations followed by the canonical source equal
    the canonical source followed by the same operations beside the fresh factors:
    \begin{align}
        T_{\mathrm{can},Y'}\bigotimes_sA_s & =\bigotimes_s(1\otimes A_s)\,T_{\mathrm{can},Y},
        \notag
    \end{align}
    where $T_{\mathrm{can},Y}$ is the operator of the canonical source of $\eta$ on
    $\{p,q\}$ from $Y$ to $Y^{pq}_{\alpha\beta}$
    \cite[Theorem~5.2]{OpenAI2026PolynomialPEPS}.
\end{theorem}
\begin{proof}\leanok
    \uses{thm:party_word_op, thm:party_add_pair_matrix, def:party_effect_matrices}
    By Theorem~\ref{thm:party_word_op} both canonical sources are preparations of
    $\eta$ along the bijections of configurations, and by
    Theorem~\ref{thm:party_add_pair_matrix} the right-hand Kronecker product is
    $1_{\alpha\times\beta}\otimes\bigotimes_sA_s$ along the same bijections. At the
    configuration $((a,b),y')$ and the input $x$ both sides have the entry
    $\eta_{ab}(\bigotimes_sA_s)_{y'x}$.
\end{proof}

\begin{theorem}[Allowed words]
    \label{thm:party_word_allowed}
    \lean{QuantumCircuit.Party.Word.isAllowed_done, QuantumCircuit.Party.Word.IsAllowed.comp,
        QuantumCircuit.Party.Word.isAllowed_localOp,
        QuantumCircuit.Party.Word.isAllowed_localOps,
        QuantumCircuit.Party.Word.isAllowed_source,
        QuantumCircuit.Party.Word.isAllowed_localPrepare,
        QuantumCircuit.Party.Word.isAllowed_localDiscard,
        QuantumCircuit.Party.Word.isAllowed_relabel,
        QuantumCircuit.Party.Word.isAllowed_sourceAt,
        QuantumCircuit.Party.Word.isAllowed_prepareProduct,
        QuantumCircuit.Party.Word.IsAllowed.toChain}
    \leanok
    \uses{def:party_word}
    The empty word, compositions of allowed words, contractions at one party and at
    every party, and relabellings are allowed words. One pair source with a normalized
    vector, in particular the canonical source of a unit vector, is allowed. The local
    preparation and the local bra of a unit vector are allowed, and so is the
    preparation of a product of unit vectors. An allowed word read as a monomial is an
    allowed monomial.
\end{theorem}
\begin{proof}\leanok
    \uses{thm:party_pair_gate_matrices, def:party_effect_matrices,
        lem:matrix_l2_contractions, thm:party_map_gates}
    A normalized source is a contraction
    (Theorem~\ref{thm:party_pair_gate_matrices}). A preparation of a unit vector is a
    contraction (Definition~\ref{def:party_effect_matrices}), and so is its adjoint;
    the column $\ket{\phi_p}$ is the preparation of $\phi_p$ on a single basis vector.
    The matrix of a bijection is a contraction (Lemma~\ref{lem:matrix_l2_contractions}).
    Reading a word as a monomial keeps allowedness by Theorem~\ref{thm:party_map_gates}.
\end{proof}

\begin{theorem}[Source lists of words]
    \label{thm:party_word_sources}
    \lean{QuantumCircuit.Party.Word.numGates_eq_sourceCount,
        QuantumCircuit.Party.Word.forallGate_card_ne_one,
        QuantumCircuit.Party.Word.gateCount_eq_sourceCount,
        QuantumCircuit.Party.Word.isNormalized_sources,
        QuantumCircuit.Party.Word.usesOnly_relabel,
        QuantumCircuit.Party.Word.UsesOnly.mem_sources,
        QuantumCircuit.Party.Word.gates_toChain}
    \leanok
    \uses{def:party_word, def:party_gate_counts}
    For every word $w$, the number of gates and the gate count are both the source
    count, since no pair source is private; the gates of $w$ read as a monomial are its
    pair sources. The source list of an allowed word is normalized. A relabelling
    that is the identification of equal bases outside $R$ uses only the parties of $R$,
    and every source of a word using only the parties of $R$ joins two parties of $R$.
\end{theorem}
\begin{proof}\leanok
    \uses{thm:party_gate_count_relations, thm:party_map_gates}
    By induction on $w$. A source has the cell $\{l,r\}$ with $l\neq r$, of two
    parties, so Theorem~\ref{thm:party_gate_count_relations} gives $\#(w)=n(w)$. The
    gates of the monomial are given by Theorem~\ref{thm:party_map_gates}. The other
    statements are read off the rules.
\end{proof}

\begin{definition}[Sources and effects of monomials]
    \label{def:party_monomial_effects}
    \lean{QuantumCircuit.Party.Monomial.sources, QuantumCircuit.Party.Monomial.effects,
        QuantumCircuit.Party.Monomial.effectParties,
        QuantumCircuit.Party.Monomial.effectCount,
        QuantumCircuit.Party.Monomial.sourceCount}
    \leanok
    \uses{def:party_pair_gates, def:party_gate_counts}
    For a monomial $M$, its sources are the data of its pair sources, the latest first,
    and its effects are the data of its pair effects in chronological order, read off
    the list of its gates. The endpoints of its effects form the list of the pairs
    $(l,r)$ of the effects. The effect count and the source count are the lengths of
    the two lists \cite[Section~5.1]{OpenAI2026PolynomialPEPS}.
\end{definition}

\begin{theorem}[Counting the sources and effects of a monomial]
    \label{thm:party_monomial_effects}
    \lean{QuantumCircuit.Party.Monomial.effectCount_add_sourceCount,
        QuantumCircuit.Party.Monomial.effectCount_comp,
        QuantumCircuit.Party.Monomial.sourceCount_comp,
        QuantumCircuit.Party.Monomial.effects_comp,
        QuantumCircuit.Party.Monomial.norm_vector_of_mem_effects,
        QuantumCircuit.Party.Monomial.effects_toChain}
    \leanok
    \uses{def:party_monomial_effects, def:party_allowed, def:party_word}
    For a monomial $M$, the effect count plus the source count is $n(M)$. Under
    composition the effect counts and the source counts add, and the effects of
    $M'\circ M$ are those of $M$ followed by those of $M'$. The effects of an allowed
    monomial are normalized, and a word read as a monomial has no effect.
\end{theorem}
\begin{proof}\leanok
    \uses{thm:party_gate_count_relations, thm:party_word_sources}
    Every gate of a monomial is a source or an effect, and the number of gates is the
    length of the list of gates (Theorem~\ref{thm:party_gate_count_relations}); count
    by induction on that list. The lists of gates of a composition are concatenated.
    Normalization follows by induction on $M$. The gates of a word read as a monomial
    are pair sources (Theorem~\ref{thm:party_word_sources}).
\end{proof}
